% =============================================================================
% BV-SALA Savings Layer — Verifiable Measurement of Cost Savings in LLM
% Applications (Open-Core Edition), v2.1, November 2026.
%
% arXiv-ready source. Compile with:
%   pdflatex main && pdflatex main
% Standard packages only (TeX Live 2023+ / arXiv default).
% The canonical source of this paper is docs/PAPER.md; this file is its
% LaTeX rendering.
% =============================================================================
\documentclass[11pt,a4paper]{article}

\usepackage[T1]{fontenc}
\usepackage[utf8]{inputenc}
\usepackage{lmodern}
\usepackage[margin=2.5cm]{geometry}
\usepackage{amsmath,amssymb,amsthm}
\usepackage{booktabs}
\usepackage{enumitem}
\usepackage{microtype}
\usepackage{hyperref}
\hypersetup{colorlinks=true,linkcolor=black,citecolor=black,urlcolor=blue}

\newtheorem{proposition}{Proposition}
\newtheorem{definition}{Definition}
\newtheorem{assumption}{Assumption}

\title{BV-SALA Savings Layer:\\ Verifiable Measurement of Cost Savings in Large Language Model Applications\\ \large (Open-Core Edition)}
\author{Joan-Marc Fisa\\ \small Founder, BiVelio Inc\\ \small \texttt{jmfisa@bivelio.com} \quad \url{https://savings.bivelio.com}}
\date{November 2026 (v2.1) \\ \small Preprint --- technical report}

\begin{document}
\maketitle

\begin{abstract}
Cost claims for large language model (LLM) optimizations are only interpretable when the
workload, the uncertainty, and the cost basis are stated explicitly. This paper evaluates
BV-SALA Savings Layer---an in-process SDK and a local gateway that combine
provider-reported usage with versioned list prices---in its \emph{open-core} distribution,
where the entire local saving stack runs with no license key, arms every local lever, and
sends zero telemetry by default. We formalize the accounting object as a function from a
request to a vector of four mutually exclusive, never-summed ledgers
(\texttt{avoided\_calls}, \texttt{eliminated\_tokens}, \texttt{compressed\_tokens},
\texttt{provider\_cached\_tokens}) and an evidence label in the lattice
\{measured, reference, estimated, unmeasured\}; we give conditions under which the cost
ratio decomposes additively over the ledgers (Proposition~1). Live double-runs under a
paired design with percentile cluster bootstrap (10{,}000 resamples, family as the
resampling unit) yield a \textbf{31.5\,\% reduction} (95\,\% CI 25.7--36.7; 100 families,
200 requests per arm) on \texttt{gpt-4o-mini}, and \textbf{75.1\,\%} (95\,\% CI
69.3--79.6; 31 families, 62 requests per arm) on \texttt{claude-sonnet-5}, of which
approximately \textbf{41 percentage points} are attributable to a prefix-cache effect
enabled by \texttt{cache\_control} in the instrumented arm rather than to a pre-existing
provider discount. A counterfactual replay of \textbf{19{,}779 gateway requests} (51
sessions, 13 days) estimates \textbf{2.49\,\%} (95\,\% CI 1.50--3.57) and is reported as
reference-valued, not measured. We add a routing-economics gradient at versioned list
prices across DeepSeek, Qwen and GLM siblings, presented as conditional arithmetic rather
than measured savings; a semantic-cache positioning against the published AWS benchmark
and GPTCache; and a preregistered roadmap specifying the factorial and non-inferiority
designs. No invoice reconciliation is claimed; every number is reproducible from public
artifacts.
\end{abstract}

\noindent\textbf{Keywords:} large language models; inference cost; provider-reported
usage; prefix caching; semantic caching; cluster bootstrap; evidence labels; open core;
model routing; measurement methodology.

\section{Introduction}\label{sec:intro}

\subsection{Problem}
The dominant cost of deploying a language model application is the provider invoice,
metered in tokens and priced by a commercial list that changes without notice. A family of
interventions---response caching, prompt compression, prompt restructuring, model routing,
prompt-cache placement---can reduce that invoice. Their literature is large
(Section~\ref{sec:related}), but the \emph{evidence} for their economic effect in a
specific deployment is usually absent: a vendor may report a percentage without a
workload, an interval, or a distinction between what the provider's counter confirmed and
what was reconstructed. This paper's thesis is that a cost claim is a measurement claim,
and inherits its obligations: an explicit instrument, a stated workload, an uncertainty
statement, and a separation between what was observed and what was inferred.

\subsection{Contributions}
\begin{enumerate}[leftmargin=1.4em]
  \item A formal accounting object (Section~\ref{sec:model}): the request$\to$report
  function, the four ledgers, the evidence lattice, and Proposition~\ref{prop:decomp}.
  \item An estimator and an uncertainty procedure (Section~\ref{sec:method}): the ratio
  estimator over paired requests with percentile cluster bootstrap at family or session
  level, and stated seeds.
  \item Live evidence (Section~\ref{sec:results}) for two providers under a paired design,
  with the Anthropic effect separated from a pre-existing provider discount, plus one
  counterfactual replay explicitly labelled as such.
  \item An open-core system contract (Section~\ref{sec:open-core}): community mode as a
  structural gate with no control plane, which removes the adoption barrier of the August
  edition while preserving every auditable guarantee.
  \item A preregistered roadmap (Section~\ref{sec:roadmap}) naming the designs this
  edition does not run.
\end{enumerate}

\subsection{Research questions}
\textbf{RQ1.} What cost reductions are observed in the specified live SDK double-runs when
provider-reported usage is valued at versioned list prices? \textbf{RQ2.} How much
uncertainty remains when dependence is handled by family- or session-level cluster
bootstrap? \textbf{RQ3.} Which components are attributable to avoided calls, to token
elimination or compression, and to cache effects enabled by the instrumented arm?
\textbf{RQ4.} What quality and latency observations are available, and which claims remain
unsupported? \textbf{RQ5.} How should a counterfactual replay and a provisional break-even
calculation be interpreted relative to live evidence?

\subsection{What changed since the August edition}
The August edition evaluated a commercial, fail-closed distribution: the layer executed
only with a verified license key, and the gateway metered Anthropic and OpenAI traffic but
not Gemini. The November release changes the system, not the measurement doctrine.
(i) \emph{Community mode}: \texttt{createSavingsLayer(\{provider\})} without a
\texttt{licenseKey} runs the full local saving stack with telemetry off, and the gateway
starts with no keys; a key that is present but invalid still fails closed.
(ii) \emph{Three protocols, one meter}: the gateway reads provider-reported usage from the
Gemini protocol, mapping \texttt{usageMetadata} onto the same usage shape as Anthropic and
OpenAI.
(iii) \emph{Chinese providers at list price}: the catalog adds DeepSeek, Qwen and GLM at
peak rates.
(iv) \emph{Zero-telemetry default}.
Nothing changes the four-ledger accounting, the evidence labels or the versioned
list-price basis.

\section{Related work}\label{sec:related}

Token pricing is a commercial decision, not a property of the Transformer. One line of
work reuses computation: PromptCache~\cite{promptcache}, CachedAttention~\cite{cachedatt},
SGLang~\cite{sglang} and CacheBlend~\cite{cacheblend} study attention or KV-cache reuse,
and Anthropic and OpenAI document provider-side prompt caching~\cite{anthropic-pc,
openai-pc}. These mechanisms can lower usage charges, but an experiment must distinguish
caching already present in both arms from caching \emph{enabled} by the intervention---the
attribution problem of Section~\ref{sec:attribution}.

A second line reduces prompt tokens: LLMLingua~\cite{llmlingua}, LLMLingua-2~\cite{llmlingua2}
and gist tokens~\cite{gist}. A third reuses responses: GPTCache~\cite{gptcache},
SCALM~\cite{scalm}, MeanCache~\cite{meancache} and vCache~\cite{vcache} address freshness,
privacy and per-prompt error bounds; P-CHR AUC and Operational Retention Rate evaluate
deployment~\cite{opgap}. Routing systems such as FrugalGPT~\cite{frugalgpt},
RouteLLM~\cite{routellm} and RouteNLP~\cite{routenlp} select among models and calibrate
escalation. Serving-efficiency work---vLLM~\cite{vllm}, SGLang~\cite{sglang}, speculative
decoding~\cite{specdec}---is orthogonal, as this paper does not measure throughput.
OpenTelemetry's GenAI conventions~\cite{otel} provide a vocabulary, not cost attribution.

BV-SALA combines established mechanisms with request-identity guards, separate ledgers,
calibrated abstention and economic gating; its differentiation is the integration of
accounting and safety, not any single mechanism. vCache is the closest comparator for
semantic caching; this work does not establish equivalent guarantees or hit rates.
Structured-data rewriting requires task- and model-specific quality evaluation: Notation
Matters shows token reductions can reduce tool-calling accuracy~\cite{notation}. The
uncertainty discussion uses established terminology~\cite{jcgm100,jcgm200} without
asserting calibration-chain traceability.

\section{System model and accounting object}\label{sec:model}

\subsection{Notation}
\begin{table}[h]
\centering
\begin{tabular}{@{}ll@{}}
\toprule
Symbol & Meaning \\
\midrule
$r$ & request: model, messages, tools, structured data, constraints \\
$u(r)=(i,o,c)$ & provider-reported usage: input, output, cache-read input tokens \\
$\pi_p=(\pi_{in},\pi_{out},\pi_{cache})$ & provider $p$'s list prices, USD per million tokens \\
$C(\pi,u)$ & cost: $(\pi_{in}i+\pi_{out}o+\pi_{cache}c)/10^6$ \\
$L$ & the layer: $r\mapsto(r',R)$ \\
$R=(C_{base}-C_{opt})/C_{base}$ & the cost ratio (primary estimand) \\
$F$ & semantic family (SDK) or session (gateway); resampling unit \\
$B$ & bootstrap resamples, $B=10{,}000$ \\
\bottomrule
\end{tabular}
\caption{Notation.}
\end{table}

All amounts are \emph{list-price} amounts at the catalog version that produced them, never
a negotiated or reconciled invoice. The internal catalog stores prices per 1{,}000 tokens;
a guard test rejects per-million values written by mistake.

\subsection{The layer}
\begin{definition}[request transform]
The layer is a (possibly partial) function $L:r\to(r',R)$ returning a transformed request
$r'$ and a report $R$. The transformation is \emph{content-preserving} under an
equivalence predicate $E(r,r')$ certified per mechanism: exact identity for cache hits, a
declared delimiter for structured-data rewriting, an output contract for constrained
decoding, header identity for upstream rewriting. Where $E$ cannot be certified, $L$ falls
back to the identity transform and reports the stage as failed-open.
\end{definition}

\begin{definition}[report]
A report is $R=(\text{calls},\text{tokens},\text{ledgers},\text{cost},\text{traces})$
where ledgers is the four-vector $(A,E,K,P)$---avoided calls, eliminated tokens,
compressed tokens, provider-cached tokens---and cost carries
$(\text{baseline},\text{actual},\text{ratio},\text{basis})$.
\end{definition}

\begin{definition}[evidence labels]
For a value $v$: \texttt{measured} (provider-reported usage plus a versioned list price),
\texttt{reference} (explicit matching against a prior \emph{executed} request whose
provider counter is reused), \texttt{estimated} (reconstructed), \texttt{unmeasured} (no
final provider counter). The labels form a partial order
$\texttt{measured}\succ\texttt{reference}\succ\texttt{estimated}\succ\texttt{unmeasured}$,
with no operation promoting a value across the order.
\end{definition}

\subsection{Ledger decomposition}\label{sec:prop}
\begin{assumption}[A1]
For a fixed request and model, the provider's usage counter is deterministic enough that
two identical requests produce the same $u$ within reported precision.
\end{assumption}
\begin{assumption}[A2]
Mechanisms act on disjoint token populations: a token avoided by call avoidance is not
also counted as eliminated; a token eliminated by rewriting is not also counted as cached.
\end{assumption}
\begin{assumption}[A3]
Prices are non-negative and the cost functional is linear in $u$.
\end{assumption}

\begin{proposition}\label{prop:decomp}
Under A1--A3, $R=R_A+R_E+R_K+R_P+\varepsilon$, where $R_A,R_E,R_K,R_P$ are the cost ratios
attributable to each ledger and $\varepsilon$ collects interaction effects between
mechanisms applied to the same request.
\end{proposition}
\begin{proof}[Proof sketch]
By A3, $C_{base}-C_{opt}=\sum_j\Delta_j$ over disjoint token populations (A2), each
$\Delta_j$ being the cost of the tokens attributable to mechanism $j$ valued at the
provider price of the token class they occupy. Dividing by $C_{base}$ yields the
decomposition; $\varepsilon$ collects the non-diagonal terms A2 does not eliminate (for
example, a compressed prompt changing the cache-placement decision).
\end{proof}

The proposition is deliberately conditional: it licenses reporting the ledgers
\emph{separately} and forbids summing them into one unconditional figure unless
$\varepsilon$ is bounded. The implementation never sums them.

\subsection{The open-core contract}\label{sec:open-core}
Community mode is a structural license gate that answers ``active, entitled to every local
feature'' with no network dependency, and an entitlements channel in the gateway grants
the four local levers without a service key. The contract: (1) with no key the layer runs
the full local stack and sends nothing; (2) with a key that is invalid, expired or revoked
the layer fails closed; (3) with a valid cloud key it additionally reports the per-request
fields enumerated in the Cloud Services Terms---never prompt content, completions or
provider credentials. Limitation: without a cloud key there is no revocation channel
(Section~\ref{sec:limits}).

\subsection{Metering the three protocols}
The gateway reads usage from Anthropic Messages, OpenAI Chat Completions/Responses and the
Gemini protocol. For Gemini, \texttt{:generateContent} and \texttt{:streamGenerateContent}
(SSE, \texttt{alt=sse} and newline-delimited JSON) are forwarded byte-for-byte; the
observer maps \texttt{usageMetadata}, discounting implicit-cache tokens from input and
summing thinking tokens into output. A stream without final \texttt{usageMetadata} is
reported \texttt{unmeasured}. These are operational checks, not savings claims.

\section{Measurement methodology}\label{sec:method}

\subsection{Design}
Each scenario runs twice against the same live provider on the same constructed content:
arm A without the layer, arm B instrumented. The unit of analysis is the semantic family---a
group of requests sharing a template and differing only in fields---so that repetition is
represented realistically. Both arms are valued with provider-reported usage at the same
catalog version.

\subsection{Estimator}
For family $f$ with $n_f$ requests per arm, $R_f=(C_{base}(f)-C_{opt}(f))/C_{base}(f)$ and
the estimand is $R=\frac1F\sum_f R_f$, treating each family as one observation: a
ratio-of-sums per family and a mean across families, not a pooled token average.

\subsection{Uncertainty}
Percentile cluster bootstrap: resample families with replacement, $B=10{,}000$, recompute
$R$, report empirical 2.5th and 97.5th percentiles. Resampling at family or session level
respects within-family dependence. Seeds and artifacts are published. This is a percentile
interval, not BCa~\cite{efron,davison}---a stated limitation.

\subsection{Attribution of the Anthropic effect}\label{sec:attribution}
The \texttt{claude-sonnet-5} reduction contains a prefix-cache effect arm B enables by
placing \texttt{cache\_control} markers. Comparing arm A's cache-read tokens (the
provider's own behavior) with arm B's isolates the enabled effect: approximately 41
percentage points of the 75.1\,\%. This is an attribution within the experiment, not a
decomposition certified by the provider.

\subsection{What the method cannot do}
It cannot establish invoice reconciliation, non-inferiority in quality or latency, or
generalization beyond the measured workloads. These are threats
(Section~\ref{sec:threats}), not buried caveats.

\section{Experimental design}\label{sec:design}
\begin{table}[h]
\centering
\small
\begin{tabular}{@{}p{3.2cm}p{5.6cm}p{5.2cm}@{}}
\toprule
Dimension & Live SDK double-run & Gateway replay \\
\midrule
Unit & semantic family (100 + 31) & session (51) \\
Requests per arm & 200 / 62 & 19{,}779 total \\
Traffic & GSM8K-derived, 30\,\% exact repeats + 20\,\% paraphrases & production sessions, one organization, 13 days \\
Models & \texttt{gpt-4o-mini}, \texttt{claude-sonnet-5} & Anthropic (Claude Code) \\
Cost basis & versioned list prices & versioned list prices, partially reconstructed baseline \\
Interval & percentile cluster bootstrap, $B=10^4$ & same at session level \\
Status & \texttt{measured} (live) & \texttt{reference} (counterfactual) \\
\bottomrule
\end{tabular}
\caption{Experimental design.}
\end{table}

\section{Results}\label{sec:results}

\begin{table}[h]
\centering
\small
\begin{tabular}{@{}llll@{}}
\toprule
Evidence, system, model & Reduction & 95\,\% CI & Workload and basis \\
\midrule
Live SDK A/B, \texttt{gpt-4o-mini} & \textbf{31.5\,\%} & [25.7; 36.7] & 100 families / 200 per arm; 30\,\% exact + 20\,\% paraphrases; SDK 0.4.31; measured \\
Live SDK A/B, \texttt{claude-sonnet-5} & \textbf{75.1\,\%} & [69.3; 79.6] & 31 families / 62 per arm; \texttt{cache\_control} in B ($\approx$41\,pp); measured \\
Counterfactual replay, Anthropic TTL & \textbf{2.49\,\%} & [1.50; 3.57] & 19{,}779 requests / 51 sessions / 13 days; partial baseline; reference \\
\bottomrule
\end{tabular}
\caption{Primary evidence (unchanged from the August edition).}
\end{table}

The OpenAI figure is the clean case: repetition-dominated traffic and a saving from calls
never executed. The Anthropic figure is large mainly because the instrumented arm enables
a provider cache the baseline did not use; both the raw figure and the enabled share are
reported. The gateway figure is small and reference-valued.

\paragraph{New operational evidence.}
Gemini protocol metering forwards \texttt{:generateContent} and
\texttt{:streamGenerateContent} byte-for-byte and maps \texttt{usageMetadata} to the same
ledger shape (non-breakage and fidelity). The DeepSeek Anthropic-format rewrite sends
\texttt{/v1/messages} to \texttt{/anthropic/v1/messages} with client headers intact
(test-locked). A catalog guard rejected a $\times$1000 per-million copy error during this
edition's development.

\subsection{Routing economics at list prices}
\begin{table}[h]
\centering
\small
\begin{tabular}{@{}lrrrl@{}}
\toprule
Model & Input & Output & Cached input & Note \\
\midrule
\texttt{claude-sonnet-5} & $\sim$2.00 & $\sim$15.00 & 0.20 & 1-h cache writes 2$\times$ \\
\texttt{gpt-4o-mini} & 0.15 & 0.60 & 0.075 & automatic caching \\
\texttt{deepseek-flash} & 0.30 (peak) & 1.20 & \textbf{0.003} & off-peak half; Anthropic-format endpoint \\
\texttt{qwen3.8-flash} & 0.15 & 0.47 & 0.015 & explicit cache, 1.25$\times$ writes \\
\texttt{glm-5.3-flash} & $\sim$0.11 & $\sim$0.39 & --- & CNY$\to$USD at 7.2 \\
\bottomrule
\end{tabular}
\caption{List prices, USD per million tokens (catalog 2026-11-03). Conditional arithmetic,
not a measurement.}
\end{table}

A cached input token is $\sim$66$\times$ cheaper on \texttt{deepseek-flash} than on
\texttt{claude-sonnet-5} ($0.20/0.003\approx 66.7$); off-peak DeepSeek input is
$\sim$13$\times$ cheaper than Sonnet input ($2.00/0.15\approx 13.3$). The cross-provider
quality gap is exactly what a cascade's gate must certify before this becomes savings; the
paper publishes the price gradient and not a routing win-rate, because no win-rate has been
measured.

\subsection{Break-even (provisional)}
At a reduction $R$ on a monthly list-price bill $B$, a seat with fee $F$ pays off iff
$RB>F$. With $F=\text{EUR~}15$ per seat per month and $R=0.315$,
$B^*=F/R\approx\text{EUR~}47.6$ per month. This is conditional arithmetic on a measured
ratio, not a forecast.

\section{Threats to validity}\label{sec:threats}
\emph{Construct:} list prices are not invoices (negotiated rates, credits, taxes,
rounding absent); every figure is labelled at list prices. \emph{Internal:} the Anthropic
effect includes a cache the instrumented arm enables, reported separately; mechanism
interactions are kept apart by Proposition~\ref{prop:decomp} with $\varepsilon$ stated.
\emph{External:} two small constructed workloads, one gateway organization, two providers.
\emph{Statistical:} percentile intervals, $B=10^4$, family-level resampling, seeds
published, no BCa claim. \emph{Reliability:} no independent replication.
\emph{Conclusion:} vendor evaluation by the product's author; every claim bounded to
reported artifacts.

\section{Limitations}\label{sec:limits}
All August limitations stand: small constructed samples (100 + 31 families), one gateway
organization, no invoice reconciliation, no independent replication, no quality or latency
non-inferiority, an in-memory linear semantic store, and the author's developer role. Two
additions: (i) the Chinese-provider rows rely on pricing pages consulted 2026-11-03, with
DeepSeek at peak rates and GLM converted at 7.2 CNY/USD; Kimi and MiniMax prices were not
verifiable and stay out of the catalog; (ii) community mode has no revocation channel by
design---the intended open-core trade, and the reason the cloud remains the paid surface.

\section{Preregistered research roadmap}\label{sec:roadmap}
\begin{enumerate}[leftmargin=1.4em]
  \item Factorial design: deduplication $\times$ \texttt{cache\_control} with per-cell power.
  \item $\rho$ sweep: exact-repeat rate swept across families (response curve).
  \item Hierarchical multi-day experiment across families, sessions, days, organizations.
  \item Non-inferiority tests: quality, freshness, latency; routing gate and TTL lever.
  \item Calibrated semantic cache at scale: multi-tenant precision, ORR~\cite{opgap}.
  \item Invoice reconciliation against one organization's provider invoices.
  \item External replication from the published branch and artifacts.
\end{enumerate}

\section{Data and code availability}
Public-bench summaries and raw CSVs: \url{https://savings.bivelio.com/benchmarks}. Software:
npm \texttt{@bivelio/savings-layer} (Apache-2.0) and PyPI
\texttt{bivelio-savings-layer}. Gateway test suites run against fake upstreams with no
network. Preprint source: \texttt{docs/PAPER.md}; LaTeX: \texttt{docs/paper/main.tex}.
Raw gateway logs remain private; aggregate replay artifacts are public.

\section{Conflict of interest}
The author is the founder of BiVelio Inc, which develops the evaluated product. No
independent replication is reported.

\begin{thebibliography}{99}
\small
\bibitem{llmlingua} H. Jiang, Q. Wu, C.-Y. Lin, Y. Yang, L. Qiu. LLMLingua: Compressing prompts for accelerated inference of large language models. EMNLP, 2023.
\bibitem{llmlingua2} Z. Pan, Q. Wu, H. Jiang, et al. LLMLingua-2: Data distillation for efficient and faithful task-agnostic prompt compression. Findings of ACL, 2024.
\bibitem{gist} J. Mu, X. L. Li, N. Goodman. Learning to compress prompts with gist tokens. NeurIPS 36, 2023.
\bibitem{promptcache} I. Gim, G. Chen, S.-s. Lee, et al. Prompt Cache: Modular attention reuse for low-latency inference. arXiv:2311.04934, 2023.
\bibitem{cachedatt} B. Gao, Z. He, P. Sharma, et al. Cost-efficient large language model serving for multi-turn conversations with CachedAttention. USENIX NSDI, 2024.
\bibitem{sglang} L. Zheng, L. Yin, Z. Xie, et al. SGLang: Efficient execution of structured language model programs. arXiv:2312.07104, 2023.
\bibitem{cacheblend} J. Yao, H. Li, Y. Liu, et al. CacheBlend: Fast large language model serving for RAG with cached knowledge fusion. EuroSys, 2025.
\bibitem{anthropic-pc} Anthropic. Prompt caching --- documentation. \url{https://docs.anthropic.com/en/docs/build-with-claude/prompt-caching}
\bibitem{openai-pc} OpenAI. Prompt caching --- API guide. \url{https://platform.openai.com/docs/guides/prompt-caching}
\bibitem{gptcache} F. Bang. GPTCache: An open-source semantic cache for LLM applications enabling faster answers and cost savings. NLP-OSS Workshop, 2023.
\bibitem{scalm} J. Li, C. Xu, F. Wang, et al. SCALM: Towards semantic caching for automated chat services with large language models. arXiv:2406.00025, 2024.
\bibitem{meancache} W. Gill, M. Elidrisi, P. Kalapatapu, et al. MeanCache: User-centric semantic caching for LLM web services. arXiv:2403.02694, 2024.
\bibitem{frugalgpt} L. Chen, M. Zaharia, J. Zou. FrugalGPT: How to use large language models while reducing cost and improving performance. arXiv:2305.05176, 2023.
\bibitem{routellm} I. Ong, A. Almahairi, V. Wu, et al. RouteLLM: Learning to route LLMs with preference data. ICLR, 2025. arXiv:2406.18665.
\bibitem{vllm} W. Kwon, Z. Li, S. Zhuang, et al. Efficient memory management for large language model serving with PagedAttention. ACM SOSP, 2023.
\bibitem{specdec} Y. Leviathan, M. Kalman, Y. Matias. Fast inference from transformers via speculative decoding. ICML, 2023.
\bibitem{otel} OpenTelemetry Project. Generative AI semantic conventions. \url{https://opentelemetry.io/docs/specs/semconv/gen-ai/}
\bibitem{jcgm100} Joint Committee for Guides in Metrology. Guide to the expression of uncertainty in measurement. JCGM 100:2008, 2008.
\bibitem{jcgm200} Joint Committee for Guides in Metrology. International vocabulary of metrology (VIM), 3rd ed. JCGM 200:2012, 2012.
\bibitem{claudecode} Anthropic. Claude Code --- agentic coding tool documentation. \url{https://www.anthropic.com/claude-code}
\bibitem{gsm8k} K. Cobbe, V. Kosaraju, M. Bavarian, et al. Training verifiers to solve math word problems. arXiv:2110.14168, 2021.
\bibitem{efron} B. Efron, R. J. Tibshirani. An Introduction to the Bootstrap. Chapman and Hall/CRC, 1993.
\bibitem{davison} A. C. Davison, D. V. Hinkley. Bootstrap Methods and Their Application. Cambridge University Press, 1997.
\bibitem{vcache} L. G. Schroeder, A. Desai, A. Cuadron, et al. vCache: Verified semantic prompt caching. ICLR, 2026. arXiv:2502.03771.
\bibitem{routenlp} D. Guo, J. Wu, S. M. Yiu. RouteNLP: Closed-loop LLM routing with conformal cascading and distillation co-optimization. arXiv:2604.23577, 2026.
\bibitem{notation} B. C. Geiger. Notation matters: A benchmark study of token-optimized formats in agentic AI systems. arXiv:2605.29676, 2026.
\bibitem{opgap} A. Baral. Closing the operational gap in semantic caching. arXiv:2606.19719, 2026.
\bibitem{bench} BiVelio. Calibrated semantic-cache evaluation and abstention report. Public benchmark documentation, September 2026. \url{https://savings.bivelio.com/benchmarks}
\end{thebibliography}

\appendix

\section{Notation and derivations}
\paragraph{A.1 Cost functional.}
$C(\pi,u)=(\pi_{in}i+\pi_{out}o+\pi_{cache}c)/10^6$. The catalog stores per-1{,}000
prices and the implementation divides by $10^3$; the two conventions produce the same
amount only when used consistently, which is why the per-million copy error the guard
rejected was material.

\paragraph{A.2 Ratio estimator.}
With per-family ratios $R_f$, $\hat R=\frac1F\sum_f R_f$ is a ratio-of-sums per family and
a mean across families. Its ratio bias is $O(1/F)$; the bootstrap interval captures
sampling variability of families, not the bias term.

\paragraph{A.3 Additivity residual.}
$\varepsilon$ is bounded in practice by the largest single-request interaction observed in
the traces; the implementation reports every stage that ran, so co-occurring mechanisms are
auditable.

\section{Bootstrap procedure}
\begin{enumerate}[leftmargin=1.4em]
  \item Compute $R_f$ from the paired arm totals of each family.
  \item Draw $B=10{,}000$ resamples of the family index set with replacement.
  \item Compute $\hat R_b$ as the mean over drawn families.
  \item Report the empirical 2.5th and 97.5th percentiles of $\{\hat R_b\}$.
  \item Publish the seed, the resample count and the artifact hash.
\end{enumerate}
The procedure is $O(BF)$; with $F\le 100$ and $B=10^4$ it runs in milliseconds and is
rerun in CI so a red test blocks a release.

\section{Reproducibility}
Install with \texttt{pnpm add @bivelio/savings-layer} (Node $\ge$ 20.19). Community mode
requires no key; the cloud key is optional. The weekly live battery and the gateway live
bench run in CI and publish each trial at \url{https://savings.bivelio.com/benchmarks}. The
catalog version that produced every figure is recorded in \texttt{src/pricing/catalogo.ts}
and echoed in each report as \texttt{versionPrecios}.

\section{Resumen en español}
Esta edición extiende la evaluación de agosto a la distribución open core: la capa corre
sin clave en modo comunitario (todas las palancas locales, telemetría cero), mide tres
protocolos de proveedor ---incluido el protocolo Gemini--- y el endpoint en formato
Anthropic de DeepSeek, y añade al catálogo versionado los precios de lista de DeepSeek,
Qwen y GLM. Las cifras primarias no cambian: 31,5\,\% [25,7; 36,7] en OpenAI y 75,1\,\%
[69,3; 79,6] en Anthropic, con $\approx$41 puntos porcentuales de efecto de caché habilitado
por el brazo instrumentado; el replay contrafactual del gateway sigue en 2,49\,\% [1,50;
3,57]. Se añade un análisis económico del enrutamiento presentado como aritmética
condicional ---no como ahorro medido---, una posición frente al benchmark publicado de
caché semántica, y un mapa de investigación prerregistrado. La contribución metodológica
es una descomposición formal del ahorro en cuatro libros que nunca se suman, con una
proposición que explicita las condiciones de aditividad y un residuo de interacción
declarado, y una jerarquía de etiquetas de evidencia que impide presentar una estimación
como un hecho de factura.

\end{document}
